\section{Real-World Experiments}
\subsection{Real-World Experimental Setup}
\begin{wrapfigure}{r}{0.55\textwidth}
    \includegraphics[width=\linewidth]{assets/figures/real_world_task_setup_fig.pdf}
    \caption{\textbf{Real-world experiment task setup}. Each task is evaluated under three conditions: In-Distribution, Task Perturbation, and Position Perturbation, to assess the policy's generalization across task and spatial variations in the real world.}
    \label{fig:real_world_exp_setup}
\end{wrapfigure}

\paragraph{Platform.} We use a Franka Emika Panda robot with two RealSense D435 cameras to  capture a front view (RGB-D) and a wrist-mounted view (RGB).

\noindent\textbf{Tasks and Evaluation.} We collect 8 pick-and-place tasks and 40 demonstrations per task for training. As shown in Fig.~\ref{fig:real_world_exp_setup}, each task is evaluated under three conditions: \textit{In-Distribution} uses the original instruction and layout; \textit{Task Perturbation} keeps the layout but swaps the target or placement region with another object, forming an unseen instruction–scene pairing; \textit{Position Perturbation} keeps the instruction but relocates the target or placement region. We run 10 rollouts per condition per task and report the average success rate. For more details, please refer to App.~\ref{supp:sec:subsec:real-world-task-setup}.

\noindent\textbf{Base Model and Baselines.} We apply our method to two pre-trained VLA backbones, $\pi_{0.5}$~\cite{pi05} and GR00T-N1.6~\cite{gr00tn1_2025}. We further compare against different representation and injection variants described in Sec.~\ref{sec:grounding_rep_inj}, all implemented using GR00T-N1.6 as the backbone.


\noindent\textbf{Grounding Source and Sub-Goal Switching.} We use Qwen3-VL-4B~\cite{qwen3vl} as the grounding source. We record the target object's z-coordinate at episode start and switch the active target from object to placement region once it rises by more than 1 cm. For more details, please refer to App.~\ref{supp:sec:subsec:real-world-grounding-source}.

\subsection{Real-World Experimental Results}
\label{sec:subsec:real-world-experimental-results}
\definecolor{mygray}{HTML}{EFEFEF}
\definecolor{myblue}{HTML}{EAF3FF}
\definecolor{deltagreen}{HTML}{3A7D44}
\definecolor{deltared}{HTML}{CC3333}

\begin{table}[t]
\centering
\caption{\textbf{Real-world experiment results.} Both backbones collapse under Task and Position Perturbation, while our method restores most of the performance and outperforms other grounding-augmented variants. The best performance in each column is highlighted in bold.}

\label{tab:real_exp}
\resizebox{0.75\linewidth}{!}{
\begin{tabular}{l ccc}
\toprule
\textbf{Method} & In Distribution & Task Perturbation & Position Perturbation \\
\midrule

% Open-Source VLA
$\pi_{0.5}$~\cite{pi05} & 85.0 & 32.5 & 27.5 \\
\rowcolor{myblue} $\pi_{0.5}$ + Ours & 91.3 & \textbf{78.8} & 61.3 \\
\rowcolor{myblue} $\Delta$ & \textcolor{deltagreen}{\footnotesize{6.3 $\uparrow$}} & \textcolor{deltagreen}{\footnotesize{46.3 $\uparrow$}} & \textcolor{deltagreen}{\footnotesize{33.8 $\uparrow$}} \\
\midrule
GR00T-N1.6~\cite{gr00tn1_2025} & 68.8 & 3.8 & 3.8 \\
\rowcolor{myblue} GR00T-N1.6 + Ours & \textbf{92.5} & 63.8 & \textbf{66.3} \\
\rowcolor{myblue} $\Delta$ & \textcolor{deltagreen}{\footnotesize{23.7 $\uparrow$}} & \textcolor{deltagreen}{\footnotesize{60.0 $\uparrow$}} & \textcolor{deltagreen}{\footnotesize{62.5 $\uparrow$}} \\
\midrule
\rowcolor{mygray} \multicolumn{4}{l}{\textit{Controlled comparison of representation and injection}} \\
\midrule
GR00T-N1.6 + 2D via Text Prompt & 87.5 & 5.0 & 7.5 \\
GR00T-N1.6 + 2D via Visual Prompt & 88.8 & 6.3 & 5.0 \\
GR00T-N1.6 + 2D via AdaLN & 87.5 & 18.8 & 16.3 \\
GR00T-N1.6 + 3D via Text Prompt & 86.3 & 5.0 & 5.0 \\
\bottomrule
\end{tabular}}
\end{table}
Tab.~\ref{tab:real_exp} reports the results on 8 real-world tasks. 
% In distribution, all grounding-augmented variants improve over the baseline, confirming the benefit of the grounding signal. Under both Task and Position Perturbation, however, the baseline and the two 2D variants collapse to near-zero success, while only our method retains most of its performance, indicating that representing grounding in 3D and injecting it directly into the action head is what enables transfer from in-distribution to OOD conditions. Notably, our real-world setup uses an off-the-shelf VLM for 2D grounding and a consumer RGB-D sensor for depth—both substantially noisier than the oracle signals in simulation—yet our method remains effective under both perturbations, confirming its practical applicability.
Our method improves both backbones under both perturbations. On $\pi_{0.5}$, success rates rise from $32.5$ and $27.5$ to $78.8$ and $61.3$ points under Task and Position Perturbation; on GR00T-N1.6, from $3.8$ and $3.8$ to $63.8$ and $66.3$. Notably, our real-world setup obtains 2D grounding from an off-the-shelf VLM and depth from a consumer RGB-D sensor---both substantially noisier than the oracle signals in simulation---yet the gains persist, confirming the practical applicability of our method.

The controlled comparison on GR00T-N1.6 reproduces the trend of Sec.~\ref{sec:grounding_rep_inj}. In distribution, every grounding-augmented variant improves over the baseline; the designs separate only under perturbation. There, both 2D prompting variants and 3D via Text Prompt collapse to near-baseline success, and even the strongest alternative, 2D via AdaLN, reaches only $18.8$ and $16.3$, far behind ours. Consistent with our analysis in Sec.~\ref{sec:grounding_rep_inj}, 2D via AdaLN still requires the policy to learn the 2D-to-3D mapping, whereas our signal already lies in 3D space, leading to better generalization. As in simulation, only the combination of a 3D representation and direct injection into the action head enables transfer from in-distribution to out-of-distribution conditions.

% Tab.~\ref{tab:libero_and_libero_pro} reports results of our method on LIBERO and LIBERO-PRO. Alongside our two controlled-comparison backbones, we include three external models as reference points: OpenVLA~\cite{kim2024openvla}, $\pi_{0}$~\cite{pi0}, and the reasoning-based MolmoAct~\cite{molmoact2025}. Although all baselines achieve strong in-distribution success, performance collapses under perturbation: OpenVLA and $\pi_{0}$ drop to $24.1$ and $32.2$ points under Task Perturbation, and to $0.0$ and $5.9$ points under Position Perturbation; even $\pi_{0.5}$ and GR00T-N1.6, trained on the same data as our method, only reach $37.3$ and $31.2$ points under Task Perturbation, and $47.0$ and $28.1$ points under Position Perturbation. This indicates that brittleness to spatial and task shifts is a widespread failure mode of current VLAs. Adding our method reverses this trend. On $\pi_{0.5}$, success rates rise to $\textbf{75.9}$ and $\textbf{72.2}$ points under Task and Position Perturbation; on GR00T-N1.6, success rates rise to $\textbf{77.5}$ and $\textbf{60.2}$ points. Consistent improvements across two backbones demonstrate that our method is backbone-agnostic. 